\chapter{Including Postscript Figures in a \LaTeX\  Document}

\centerline[This chapter is contributed by Joan Carletta]

This chapter describes how to include a postscript file in a \LaTeX\
document. This capability allows \LaTeX\ users to include figures,
drawings, and graphs more easily generated using software tools other
than \LaTeX\ in their documents. For example, a document may contain:
\begin {itemize}
  \item graphs generated with ``\verb+gnuplot+''.\\
        ``\verb+gnuplot+'' can generate postscript format directly.
  \item figures drawn with ``\verb+xfig+''.\\
        ``\verb+xfig+'' generates files with a ``\verb+.fig+'' extension.  
        This can be translated into postscript format by using
        the ``\verb+fig2dev+'' utility.  For best results, use the
        command with the \verb+-P+ and \verb+-c+ options:
        \begin {quotation}
          \verb+fig2dev -Lps -P -c +{\em figure}\verb+.fig > +{\em figure}\verb+.ps+
        \end {quotation}
  \item X window dumps generated by ``\verb+xwd+''.\\
        ``\verb+xwd+'' generates files in a special X window dump
        format.  Suppose that you have dumped an image of an X window 
        in the file ``{\em figure}'' by typing 
        \begin {quotation}
          \verb+xwd > +{\em figure}
        \end {quotation}
        in some other window, and then clicking the rightmost mouse
        button on the window that you'd like to capture.
        To convert this file into postscript format, use the
        ``\verb+xpr+'' utility:
        \begin {quotation}
          \verb+xpr -device ps -output +{\em figure}\verb+.ps -portrait -compact +{\em figure}
        \end {quotation}
\end {itemize}
Many other tools can also be used to produce files in the postscript
format.

Figure \ref{figure1} outlines the steps taken to include a postscript
file. For this example, the postscript file is included as a figure
of a \LaTeX\ document. The figure that results from the code of
Figure \ref{figure1} is shown in Figure \ref{this_figure}.

\begin {figure}[p]
  \begin {center}
%   \fbox {
 \begin {minipage}{4.8in}
\begin{verbatim}
  \begin {document}
\end{verbatim}
   \hspace {.5in} $\vdots$
\begin{verbatim}
    \input psfig
\end{verbatim}
   \hspace {.5in} $\vdots$
\begin{verbatim}
    \begin{figure}
      \vskip0.2truein
      \centerline {\hbox{\psfig {file=figure.ps,height=2.0in}}}
      \vskip0.1truein
      \caption {This is a figure.}
      \label{this_figure}
    \end{figure}
\end{verbatim}
   \hspace {.5in} $\vdots$
\begin{verbatim}
  \end{document}
\end{verbatim}
  \end {minipage}
%}
  \end {center}
  \caption {Steps for including a postscript file as part of a figure.}
  \label{figure1}
\end {figure}

The ``\verb+\input psfig+'' line reads in the file
``\verb+psfig.tex+'', which defines the command ``\verb+\psfig+'' that
is used to include postscript figures.  The basic command to include a
figure is:
\begin {quotation}
\verb+\psfig {file=+{\em figure}\verb+.ps,height=2.0in}+
\end {quotation}
Everything else has been added for cosmetic reasons. The argument to
the ``\verb+\psfig+'' command has several components, separated by
commas; here, the ``\verb+file=+'' component specifies the name of the
postscript file to be included, and the ``\verb+height=+'' component
specifies the height of the box that \LaTeX\ should create for the
figure. The figure will be scaled to fit within a box of the specified
height. Alternatively, the width of the box may be specified with a
``\verb+width=+'' component. It is not a good idea to specify both
height and width, since \LaTeX\ may stretch the figure unnaturally in
one direction to meet both constraints. Other options are available,
and are documented at the beginning of the ``\verb+psfig.tex+'' file,
which is in the ``\verb+/usr/local/lib/tex/inputs+'' directory.

\begin{figure}[p]
  \vskip0.2truein
  \centerline {\hbox{\psfig {file=figure-sec.ps,height=2.0in}}}
  \vskip0.1truein
  \caption {This is a figure.}
  \label{this_figure}
\end{figure}


